\section{Experimental evidence and parameter identification}\label{sec:experiments}

A wave theory is only as useful as its parameters. The section asks which experiment probes which parameter, what the experiments can and cannot discriminate, and how published lag times and fractional orders should be read.

\subsection{Second sound and ballistic heat pulses}
The direct signature of wave-like heat transport is second sound, a propagating temperature wave. It was predicted for dielectric crystals at low temperature \citep{Chester1963,GuyerKrumhansl1966} and observed in solid helium \citep{Ackerman1966} and in sodium fluoride \citep{McNelly1970,Jackson1970}, in narrow temperature windows where momentum-conserving phonon collisions dominate over resistive ones. A further observation in graphite above $100$\,K \citep{Huberman2019} extends the window to a much higher temperature in a material of high anisotropy. There the fitted second-sound speed at $85$\,K is about $3200$\,m/s against $3650$\,m/s from ab initio calculations, with oscillations at $104$ and $125$\,K and none at $50$ and $150$\,K (values from the first-release version of the paper, with no error bar quoted for the fit). These experiments work in a hydrodynamic phonon window with $\tau_N\ll\tau_R$, exactly the conditions under which the Guyer--Krumhansl law (\cref{sec:kinetic}) applies. They do not show that a Cattaneo or lag law with a single relaxation time describes conduction in ordinary solids at room temperature. A recent review of non-Fourier phonon conduction at the micro- and nanoscale \citep{Chen2021NRP} sorts the observations into Casimir--Knudsen, hydrodynamic (second-sound), coherent-localization and low-dimensional divergence regimes and assesses the experimental evidence for each.

\subsection{Ultrafast optical methods}
Pump--probe reflectivity measures electron--lattice energy exchange on femtosecond to picosecond timescales \citep{Brorson1987} and picosecond ultrasonics generates and detects coherent acoustic pulses in thin films \citep{Thomsen1986}. The first is where an electron--lattice two-temperature model (distinct from the conductive and thermodynamic temperatures of the two-temperature theories of \cref{sec:lags}) provides the continuum model, and what is measured is an electron--phonon coupling time rather than $\tau_q$. In the second, the heat input launches a strain pulse whose amplitude and shape depend on diffusion within the first tens of nanometers, so thermoelastic coupling acts as the source and not as a loss mechanism. Time-domain thermoreflectance \citep{Cahill2003} fits a diffusive model with interface resistance to extract the conductivity of films and interfaces, and a short-delay deviation from that fit is where lag or ballistic corrections are usually claimed.

\subsection{Resonator quality factors}
Micro- and nano-resonators provide the cleanest tests of thermoelastic damping because $Q$ can be measured against temperature and thickness. The quantitative signature, \cref{eq:zenerQ,eq:LR}, is a loss peak at $f_Z=\pi\chi/(2h^2)$ with the scaling $f_Z\propto h^{-2}$ and a height proportional to $\Delta$. Verification therefore needs a sweep of thickness or temperature rather than a single value, with anchor and surface losses subtracted; those losses lie outside the models of this review. Only after the contributions of the lag and the size effect have been bounded can a mismatch be attributed to either of them.

\subsection{Heterogeneous media and biological tissue}
The relaxation times of order seconds reported for non-homogeneous materials, processed biological tissue among them \citep{Mitra1995,Kaminski1990}, reflect heat transfer through a heterogeneous microstructure and are effective macroscopic parameters. Reading them as molecular relaxation times is not warranted, and the values have been disputed. For tissue the parabolic Pennes bioheat equation \citep{Pennes1948} remains the baseline against which any hyperbolic correction has to be judged.

\subsection{What the dispersion relation can identify}\label{sec:identify}
Equation \eqref{eq:disp} contains the heat law only through the single function $G(s)$ and is \emph{linear} in $G$. If the complex wavenumber $k(\omega)$ of the elastic-type branch is measured (phase velocity from $\operatorname{Re}k$, attenuation from $\operatorname{Im}k$) and $\rho,c_1,c,\beta,T_0$ are known, \eqref{eq:disp} can be solved for $G$ at each frequency,
\begin{equation}
G(-\ii\omega)=-\,\frac{s\,\bigl[\rho^2c\,(c_1^2K-\omega^2)+\beta^2T_0K\bigr]}{\rho\,K\,(c_1^2K-\omega^2)},\qquad K=k^2,\ s=-\ii\omega .\label{eq:Ginvert}
\end{equation}
A measured $k(\omega)$ therefore determines the heat-flux function $G(-\ii\omega)$ \emph{without assuming} a lag model; the model enters afterwards, when a parametric form of $G$ is fitted.

For the template $G(s)=[k(1+\tau_Ts)+k^*(1+\tau_vs)/s]/(1+\tau_qs)$ the five parameters $(k,k^*,\tau_q,\tau_T,\tau_v)$ are \emph{not} all identifiable. Writing the numerator as $k^*+(k+k^*\tau_v)s+k\tau_Ts^2$, $G$ depends only on the four combinations
\begin{equation}
k^*,\qquad k+k^*\tau_v,\qquad k\,\tau_T,\qquad \tau_q ,
\end{equation}
so the parameters can move along $(\delta k,\delta\tau_T,\delta\tau_v)=(\delta,\,-\tau_T\delta/k,\,-\delta/k^*)$ without changing $G$ and hence without changing any wave property of the one-dimensional problem, which is consistent with the combination $\tau_v^{*}=k+k^*\tau_v$ of Roy Choudhuri \citep{RoyChoudhuri2007}. Two consequences follow: dispersion or attenuation data alone cannot separate $k$ from $\tau_v$, and a pair $(k,\tau_v)$ fitted to such data is a convention rather than a measurement (compare \cref{sec:delayfid}).

\paragraph{Numerical check (synthetic data).} The identifiability check generates $k(\omega)$ from the template with a chosen parameter set, verifies that \eqref{eq:Ginvert} returns $G$ to $10^{-9}$, that the Jacobian of $G(-\ii\omega)$ with respect to the five parameters has numerical rank four (singular values relative to the largest: $1,\ 0.93,\ 0.44,\ 0.05,\ 3\times10^{-11}$), and that two parameter sets differing along the null direction give identical $G$. With $0.5\,\%$ random noise on complex $k$ at twelve frequencies, the combinations $k^*$ and $\tau_q$ are recovered to $2\,\%$ and below $1\,\%$, $k+k^*\tau_v$ to about $12\,\%$ and $k\tau_T$ only to about $26\,\%$ (\cref{fig:identify}). These figures describe this synthetic example, with one parameter set, one noise level and one seed, and not any material; no experimental data enter.

\begin{figure}[t]
\centering
\includegraphics[width=\textwidth]{figs/fig_identify.pdf}
\caption{Synthetic identification test. Left: $G(-\ii\omega)/k$ recovered from $0.5\,\%$-noisy complex wavenumbers through \eqref{eq:Ginvert} (markers) against the exact function (lines). Right: singular values of the Jacobian of $G$ with respect to $(k,k^*,\tau_q,\tau_T,\tau_v)$; the fifth is zero to rounding, the fourth is small.}\label{fig:identify}
\end{figure}


\subsection{How to read a published lag time}
\begin{enumerate}
\item Identify what was measured: a response time of a temperature signal, a conductivity, or a quality factor.
\item Identify the model used to extract $\tau$ and whether other parameters were fitted simultaneously. A fit of two or more lag parameters from a single curve is usually not identifiable (compare \cref{sec:delayfid}).
\item Compare $\omega\tau$ with the validity limits of the model (\cref{sec:delayfid} for truncated lag laws, \cref{sec:kinetic} for the kinetic basis).
\item Check whether a simpler model (Fourier with interface resistance, or Guyer--Krumhansl) explains the same data.
\end{enumerate}
A model parameter that cannot be measured independently of the model that defines it is a fitting constant and not a material property. Papers that tabulate such constants for many materials without an independent test provide limited guidance.

\begin{summarybox}[Summary]
Experimental evidence: second sound at cryogenic temperatures (Ackerman--Bertman, Jackson et al.) is the cleanest demonstration of wave-like heat transport. Room-temperature reports of large relaxation times in nonhomogeneous materials are debated. Pump--probe methods measure electron--phonon energy exchange, and picosecond ultrasonics uses thermoelastic coupling as a source; neither yields a lag time without a model fit. A model parameter not measurable independently of the model is a fitting constant, not a material property.
\end{summarybox}
